\chapter{形状与插值强度的性质}
\label{ch:parameters}
\section{路径、覆盖与局部速度}
固定 $q,T$ 时，$\eta$ 从零增加到一，把目标教师从锚点移向教师。对固定 $\kappa$，任意凸生成散度 $D_\varphi(m_{\kappa,\eta}\Vert q)$ 随 $\eta$ 非减；不同 $\kappa$ 下，相同 $\eta$ 不代表相同距离。其交点证明见\appref{app:shape}。

正形状有一个分布无关的覆盖保证：
\begin{equation}
m_a\ge(1-\eta)^{1/\kappa}q_a+\eta^{1/\kappa}T_a,\qquad 0<\kappa\le1.
\label{eq:coverage-main}
\end{equation}
证明使用 $h_a\le(1-\eta)q_a+\eta T_a$，故 $Z\le1$；以及 $(x+y)^{1/\kappa}\ge x^{1/\kappa}+y^{1/\kappa}$。于是 $m_a=h_a/Z\ge h_a$ 给出下界。几何端点没有统一的正教师覆盖系数。

对固定严格正分布，小步展开为
\begin{equation}
m_a=q_a+\eta v_{\kappa,a}+O(\eta^2),\qquad
v_{\kappa,a}=\frac{q_a}{\kappa}
\left[(T_a/q_a)^\kappa-\E_q(T/q)^\kappa\right].
\label{eq:velocity-main}
\end{equation}
$\kappa=0$ 时取连续极限 $v_{0,a}=q_aC_q\log(T_a/q_a)$。这是 $D_\kappa(T\Vert q)$ 在 Fisher 度量下的负自然梯度。形状改变高维方向，不能通常由一个学习率缩放替代。

\section{监督尾部与方差阈值}
当 $\kappa>0,\eta<1$ 时，
\begin{equation}
A_{\kappa,\eta}(x)\ge-c_{\kappa,\eta},\qquad
c_{\kappa,\eta}=-\frac{\log(1-\eta)}{\kappa\eta}.
\end{equation}
正尾部仍无逐点统一上界，但在 $q$ 下利用 $\E_q(T/q)=1$ 可以控制矩。令 $V_{\kappa,\eta}=\Var_q A_{\kappa,\eta}(T/q)$。对每个固定正 $\eta<1$ 和 $\kappa>0$，$V$ 在所有有限正概率分布上有统一有限上界。

更强的小步问题需要另一个结论：
\begin{equation}
\sup_{q,T}V_{\kappa,\eta}
=\Theta_\kappa\!\left(\eta^{\min\{2,1/\kappa\}-2}\right),
\qquad \eta\downarrow0,\quad\kappa>0.
\label{eq:variance-threshold-main}
\end{equation}
所以保持统一小步方差有界的区间为 $0<\kappa\le1/2$。几何端点对任意固定正 $\eta$ 都没有分布无关界。这里上确界允许分布随 $\eta$ 改变；固定一对正分布时，每种形状的小步极限都有限。

对 $0<\kappa\le1/2$，还有误差相关界
\begin{equation}
V_{\kappa,\eta}
\le\frac{D_\kappa(T\Vert q)}{\kappa(1-\eta)^2}
\le\frac{K_T}{\kappa(1-\kappa)(1-\eta)^2}.
\label{eq:error-variance-main}
\end{equation}
若 score 范数统一不超过 $S$，中心化教师梯度的协方差迹不超过 $S^2V$。标量监督方差与完整参数梯度方差由这个额外条件连接；精确均值中心本身不保证 score 梯度的方差最小。

\section{平方根点的作用与限制}
令 $a=\sum_v\sqrt{q_vT_v}$、$w=1-\eta$。平方根目标满足
\begin{equation}
\sqrt m=\frac{w\sqrt q+\eta\sqrt T}{\sqrt Z},\qquad
Z=w^2+\eta^2+2w\eta a.
\end{equation}
因为 $Z\le1$，
\begin{equation}
\sum_v\sqrt{m_vT_v}=\frac{wa+\eta}{\sqrt Z}\ge wa+\eta,
\quad
H^2(m,T)\le(1-\eta)H^2(q,T).
\label{eq:hellinger-contraction-main}
\end{equation}
这是精确分布更新的全局收缩。所有形状在固定正教师附近却共享
\begin{equation}
m_{\kappa,\eta}(T+\delta,T)
=T+(1-\eta)\delta+O(\|\delta\|^2).
\label{eq:local-rate-main}
\end{equation}
因此平方根的全局 Hellinger 保证不代表更好的局部一阶概率收缩率。

平方根点唯一最小化小步监督方差相对于自身教师散度的最坏比值；若改为最小化绝对最坏方差，最优形状约为 $0.441928$。两种准则的推导在\appref{app:shape}。固定 $q,\eta$ 时，BPO 的独立 logits Hessian 为
\begin{equation}
\nabla_z^2L_{\rm BPO}=\eta^{-1}[\diag(q)-qq^\top],
\end{equation}
其中不含 $\kappa$。形状改变目标和初始残差，不能由此宣称改善该二次问题的条件数。

\section{共同收敛界与形状比较}
对同一正教师、$0\le\kappa\le1$ 和 $0<\eta<1$，精确目标共享
\begin{equation}
\KL(T\Vert m_{\kappa,\eta})\le(1-\eta)\KL(T\Vert q).
\label{eq:forward-kl-main}
\end{equation}
证明使用未归一化坐标 $q_a^{1-\eta}T_a^\eta\le h_a\le(1-\eta)q_a+\eta T_a$ 及 $Z\le1$。因此精确外层有全局几何上界；教师附近的双向 KL 比例则共同趋于 $(1-\eta)^2$。完整证明、双锚点扰动界及反例见\appref{sec:shape-comparison}。

\begin{center}
\begin{tabular}{lp{36mm}p{39mm}p{36mm}}
\toprule
形状 & 精确更新的坐标 & 教师概率覆盖 & 小步标准化方差\\
\midrule
$0$ & log-odds 线性插值 & 无统一正覆盖系数 & 无统一有限上界\\
$1/2$ & 平方根球面插值 & 至少 $\eta^2T_a$ & 最坏极限为 $4$\\
$1$ & 概率线性插值 & 至少 $\eta T_a$ & 最坏值按 $1/\eta$ 增长\\
\bottomrule
\end{tabular}
\end{center}

算术端点对所有凸生成散度具有 $(1-\eta)$ 的乘法收缩，平方根的主要优势是小步矩控制及相应的误差相关界。固定 $\eta$ 的更新后教师 KL 可以由不同端点取得最小值，因此不存在逐实例的统一形状排序。固定强度、固定 Fisher 移动速度和固定参数更新预算是三种不同的比较口径。

\section{外推的适用范围}
$\eta>1$ 把插值改为外推。几何端点对任意正 $\eta$ 保持正性；正形状要求 $(T_a/q_a)^\kappa>1-1/\eta$ 对全部动作成立。对固定正教师，算术外推在 $\eta=2$ 形成合法两周期；平方根在同一强度下，只要第一步严格为正，精确迭代仍收敛，双向 KL 渐近为 $1/n$。更一般的 $0<\kappa<1$ 在教师附近有 $[4\kappa(1-\kappa)n]^{-1}$ 的边界 KL 率。

这些结论不改变正文的默认插值范围，也不把精确边界收敛解释为采样训练的最佳强度。平方根在 $1<\eta<2$ 的输出扰动稳态界含 $1/(2-\eta)$，接近边界时会变差。正性、角度递推、两周期和完整高阶证明见\appref{sec:positive-extrapolation}。
